% =====================================================================
%  Taller de tipos de requisitos (ICR_QC_01 / TL01)
%  Compilar con pdfLaTeX + Biber (configuración predeterminada de Overleaf)
% =====================================================================
\documentclass[11pt,a4paper]{article}

% ---------- Codificación e idioma ----------
\usepackage[utf8]{inputenc}
\usepackage[T1]{fontenc}
\usepackage{lmodern}
\usepackage[spanish,es-tabla]{babel}

% ---------- Página ----------
\usepackage[a4paper,left=2.5cm,right=2.5cm,top=2.5cm,bottom=2.5cm]{geometry}
\usepackage{microtype}
\usepackage{setspace}
\setstretch{1.12}
\setlength{\parskip}{0.5em}
\setlength{\parindent}{0pt}

% ---------- Tablas, figuras y color ----------
\usepackage{xcolor}
\usepackage{array}
\usepackage{colortbl}
\usepackage{longtable}
\usepackage{ragged2e}
\usepackage{graphicx}
\usepackage{float}
\usepackage{enumitem}
\usepackage{needspace}
\usepackage{pdflscape}
\usepackage[font=small,labelfont=bf,labelsep=period,skip=6pt]{caption}

% ---------- Títulos y encabezados ----------
\usepackage{titlesec}
\titleformat{\section}{\Large\bfseries}{\thesection.}{0.6em}{}
\titleformat{\subsection}{\large\bfseries}{\thesubsection.}{0.6em}{}
\titlespacing*{\section}{0pt}{1.6em}{0.8em}
\titlespacing*{\subsection}{0pt}{1.2em}{0.6em}

\usepackage{fancyhdr}
\pagestyle{fancy}
\fancyhf{}
\fancyhead[L]{\small\itshape Taller de tipos de requisitos}
\fancyhead[R]{\small\itshape ICR\_QC\_01}
\fancyfoot[C]{\small\thepage}
\renewcommand{\headrulewidth}{0.4pt}
\renewcommand{\footrulewidth}{0pt}

% ---------- Bibliografía (citas numéricas estilo IEEE) ----------
\usepackage{csquotes}
\usepackage[backend=biber,style=ieee,sorting=none]{biblatex}
\addbibresource{referencias.bib}
\DefineBibliographyStrings{spanish}{url={Disponible en}}

% ---------- Hipervínculos ----------
\usepackage[colorlinks=true,linkcolor=black,citecolor=blue!50!black,urlcolor=blue!50!black,
            pdftitle={Taller de tipos de requisitos},
            pdfauthor={},
            pdfsubject={Ingeniería de Requerimientos}]{hyperref}
\usepackage{xurl}

% ---------- Formato de tablas ----------
\setlength{\tabcolsep}{3pt}
\renewcommand{\arraystretch}{1.25}
\newlength{\tw} % ancho útil de las columnas de una tabla
\newcolumntype{L}[1]{>{\RaggedRight\arraybackslash}p{#1}}
\newcommand{\encabezado}[1]{\textbf{#1}}
\newcommand{\filaencabezado}{\rowcolor{gray!22}}

% ---------- Marcador de imágenes (se sustituye al añadir el archivo) ----------
% Uso: \capturaimagen{nombre-sin-extension}{texto del marcador}
\newcommand{\capturaimagen}[2]{%
  \IfFileExists{imagenes/#1.png}{\includegraphics[width=0.9\linewidth,height=0.45\textheight,keepaspectratio]{imagenes/#1.png}}{%
  \IfFileExists{imagenes/#1.jpg}{\includegraphics[width=0.9\linewidth,height=0.45\textheight,keepaspectratio]{imagenes/#1.jpg}}{%
  \IfFileExists{imagenes/#1.jpeg}{\includegraphics[width=0.9\linewidth,height=0.45\textheight,keepaspectratio]{imagenes/#1.jpeg}}{%
  \IfFileExists{imagenes/#1.pdf}{\includegraphics[width=0.9\linewidth,height=0.45\textheight,keepaspectratio]{imagenes/#1.pdf}}{%
  \fcolorbox{gray}{gray!8}{\parbox[c][5.5cm][c]{0.85\linewidth}{\centering\small\itshape #2}}}}}}}

% ---------- Datos de la portada (editar aquí) ----------
\newcommand{\datoUniversidad}{Universidad Técnica Estatal de Quevedo (UTEQ)}
\newcommand{\datoFacultad}{Facultad de Ciencias de la Computación}
\newcommand{\datoCarrera}{Ingeniería en Software}
\newcommand{\datoAsignatura}{Ingeniería de Requerimientos}
\newcommand{\datoCodigoActividad}{ICR\_QC\_01}
\newcommand{\datoFormato}{TL01}
\newcommand{\datoTipoActividad}{Formativa}
\newcommand{\datoTaller}{Taller de tipos de requisitos}
\newcommand{\datoTitulo}{Identificación y clasificación de tipos de requisitos}
\newcommand{\datoEstudiante}{Rivera Rodriguez Joseph Jeremy}
\newcommand{\datoNivel}{Quinto semestre}
\newcommand{\datoParalelo}{A}
\newcommand{\datoDocente}{GUERRERO ULLOA GLEISTON CICERON, PhD}
\newcommand{\datoFecha}{7 de noviembre de 2026}

\begin{document}

% ---------- Portada (sin numeración) ----------
\pagenumbering{gobble}
% Portada académica. Los datos se editan en el preámbulo de main.tex (comandos \dato...).
\begin{titlepage}
\thispagestyle{empty}
\centering
\setstretch{1.1}

% Para añadir el escudo/logo de la universidad, coloque el archivo en imagenes/ y
% descomente la siguiente línea:
% \includegraphics[height=3cm]{imagenes/logo-utq.png}\\[1.2cm]

\vspace*{1.2cm}

{\huge\bfseries \datoUniversidad\par}
\vspace{0.9cm}
{\Large\bfseries \datoFacultad\par}
\vspace{0.5cm}
{\large Carrera de \datoCarrera\par}

\vspace{1.0cm}
\rule{0.75\linewidth}{0.8pt}\par
\vspace{0.9cm}

{\large\scshape \datoTaller\par}
\vspace{0.5cm}
{\LARGE\bfseries \datoTitulo\par}

\vspace{0.9cm}
\rule{0.75\linewidth}{0.8pt}\par
\vspace{1.1cm}

{\renewcommand{\arraystretch}{1.55}
\small
\begin{tabular}{@{}r@{\hspace{0.8em}}l@{}}
  \textbf{Asignatura:}           & \datoAsignatura \\
  \textbf{Tipo de actividad:}    & \datoTipoActividad \\
  \textbf{Estudiante:}           & \datoEstudiante \\
  \textbf{Nivel:}                & \datoNivel \\
  \textbf{Paralelo:}             & \datoParalelo \\
  \textbf{Docente:}              & \datoDocente \\
\end{tabular}\par}

\vfill
{\normalsize\textbf{Fecha de entrega:} \datoFecha\par}
\vspace{0.5cm}
\end{titlepage}

\clearpage
\pagenumbering{arabic}
\setcounter{page}{1}
\pagestyle{fancy}

% ---------- Contenido ----------
\section{Clasificación de los requisitos}

Para realizar la clasificación se consideran tres aspectos: el tipo de requisito, su relación con el producto o el proceso y el nivel al que se aplica, es decir, al sistema completo o específicamente al software.

La ISO/IEC/IEEE 29148:2018 establece criterios para la ingeniería de requisitos y contempla la especificación de necesidades, funciones, restricciones y condiciones que debe cumplir una solución. También proporciona criterios para que los requisitos estén redactados de manera clara, precisa y verificable~\cite{iso29148}.

Por otro lado, la ISO/IEC 25010:2023 presenta un modelo de calidad del producto compuesto por nueve características y sus respectivas subcaracterísticas. Este modelo permite identificar y evaluar las propiedades de calidad que debe cumplir un producto de software o un sistema~\cite{iso25010}.

\begingroup
\footnotesize
\setlength{\tw}{\dimexpr\linewidth-10\tabcolsep-6\arrayrulewidth\relax}
\begin{longtable}{|L{0.08\tw}|L{0.23\tw}|L{0.12\tw}|L{0.12\tw}|L{0.45\tw}|}
\caption{Clasificación de los requisitos E01--E20.}\label{tab:clasificacion}\\
\hline
\filaencabezado \encabezado{Código} & \encabezado{Tipo de requisito y característica ISO/IEC 25010:2023} & \encabezado{Producto / proceso} & \encabezado{Sistema / software} & \encabezado{Justificación} \\ \hline
\endfirsthead
\multicolumn{5}{@{}l}{\footnotesize\itshape Continuación de la Tabla~\thetable} \\ \hline
\filaencabezado \encabezado{Código} & \encabezado{Tipo de requisito y característica ISO/IEC 25010:2023} & \encabezado{Producto / proceso} & \encabezado{Sistema / software} & \encabezado{Justificación} \\ \hline
\endhead
\hline \multicolumn{5}{r}{\footnotesize\itshape Continúa en la página siguiente} \\
\endfoot
\hline
\endlastfoot

E01 & Funcional & Producto & Software & Establece la operación que permite al paciente reservar una cita seleccionando especialidad, médico, fecha y hora. Describe una función que debe ofrecer la aplicación~\cite{iso29148}. \\ \hline
E02 & Funcional & Producto & Software & Especifica el envío automático de un recordatorio por SMS antes de la cita. Se trata de una función que debe ejecutar la plataforma~\cite{iso29148}. \\ \hline
E03 & Calidad: eficiencia de desempeño – comportamiento temporal & Producto & Software & Se refiere al tiempo que tarda la aplicación en responder. El comportamiento temporal contempla los tiempos de respuesta y las tasas de procesamiento que debe cumplir el producto bajo determinadas condiciones~\cite{iso25010}. \\ \hline
E04 & Funcional & Producto & Software & Define la posibilidad de registrar la llegada de un paciente desde recepción. Es una operación concreta que debe permitir el sistema~\cite{iso29148}. \\ \hline
E05 & Restricción tecnológica & Producto & Software & Obliga a utilizar Java 21 y PostgreSQL, por lo que limita las tecnologías que pueden seleccionarse para implementar el software. Se considera una condición impuesta a la solución~\cite{iso29148}. \\ \hline
E06 & Funcional & Producto & Software & Establece que el médico puede consultar los pacientes programados para el día y registrar si fueron atendidos. Describe funciones que debe ofrecer la aplicación~\cite{iso29148}. \\ \hline
E07 & Calidad: seguridad – confidencialidad e integridad & Producto & Software & Se refiere a la protección de los datos médicos. La confidencialidad busca evitar el acceso a información por personas no autorizadas, mientras que la integridad protege los datos frente a modificaciones no autorizadas~\cite{iso25010}. \\ \hline
E08 & Restricción de compatibilidad & Producto & Software & Establece Android 10 como versión mínima para ejecutar la aplicación móvil. Es una condición técnica que limita el entorno en el que podrá funcionar el software~\cite{iso29148}. \\ \hline
E09 & Funcional & Producto & Software & Define cómo debe registrarse una cancelación realizada con menos de cuatro horas de anticipación. Es una regla de funcionamiento que debe aplicar la plataforma~\cite{iso29148}. \\ \hline
E10 & Calidad: capacidad de interacción – inclusividad y operabilidad & Producto & Software & Se relaciona con la interacción de los pacientes adultos mayores con la aplicación. La inclusividad considera las necesidades de diferentes usuarios, mientras que la operabilidad se refiere a la facilidad para utilizar y controlar las funciones del producto~\cite{iso25010}. \\ \hline
E11 & Restricción de plazo & Proceso & Sistema & Establece un tiempo máximo de seis meses para completar el proyecto. Limita la planificación y ejecución del proceso de desarrollo~\cite{iso29148}. \\ \hline
E12 & Funcional & Producto & Software & Especifica que la plataforma debe cerrar automáticamente la sesión después de diez minutos de inactividad. Aunque contribuye a la seguridad, el enunciado describe una acción concreta que debe realizar el software~\cite{iso29148}. \\ \hline
E13 & Funcional & Producto & Sistema & Define el funcionamiento del quiosco: recibir el número de identificación mediante una pantalla táctil y emitir un ticket. Se considera de sistema porque intervienen tanto el software como los dispositivos físicos~\cite{iso29148}. \\ \hline
E14 & Calidad: fiabilidad – disponibilidad & Producto & Software & Establece que la plataforma debe estar operativa cuando los usuarios la necesiten. La disponibilidad se relaciona con la capacidad del producto de estar operativo y accesible cuando se requiere~\cite{iso25010}. \\ \hline
E15 & Funcional & Producto & Software & Permite al director médico generar reportes mensuales de inasistencias por especialidad. Es una función específica que debe ofrecer la plataforma~\cite{iso29148}. \\ \hline
E16 & Restricción del proceso & Proceso & Sistema & Establece que cada incremento debe ser validado con el director médico al finalizar una iteración de dos semanas. Determina una condición que debe seguirse durante el proceso de desarrollo~\cite{iso29148}. \\ \hline
E17 & Restricción legal & Producto & Sistema & Exige cumplir la Ley Orgánica de Protección de Datos Personales del Ecuador. Es una obligación que condiciona el tratamiento de la información y el funcionamiento del sistema~\cite{iso29148}. \\ \hline
E18 & Calidad: compatibilidad – interoperabilidad & Producto & Software & Establece que la plataforma debe intercambiar información con el sistema de facturación existente. La interoperabilidad se relaciona con la capacidad de intercambiar información y utilizarla entre diferentes productos~\cite{iso25010}. \\ \hline
E19 & Calidad: fiabilidad – recuperabilidad & Producto & Sistema & Establece condiciones para recuperar el servicio después de una falla del servidor, sin perder las citas confirmadas. La recuperabilidad contempla restaurar los datos afectados y el estado esperado del sistema~\cite{iso25010}. \\ \hline
E20 & Restricción del proceso & Proceso & Software & Obliga a que los cambios en el código sean revisados por otro integrante antes de incorporarse a la rama principal. Es una condición del procedimiento de desarrollo y control del software~\cite{iso29148}. \\ \hline
\end{longtable}
\endgroup



\section{Identificación y reformulación de requisitos de calidad}

Los requisitos de calidad seleccionados contienen expresiones generales que no permiten comprobar objetivamente su cumplimiento. Algunas palabras utilizadas, como «rápida», «seguros», «fácil», «siempre» y «sin problemas», no establecen condiciones específicas para determinar si el sistema cumple con lo solicitado.

Según la ISO/IEC/IEEE 29148:2018, los requisitos deben expresar las necesidades de manera clara, precisa y sin ambigüedades, de forma que puedan ser evaluados y verificados~\cite{iso29148}. Por este motivo, es necesario reformular los enunciados seleccionados, incorporando métricas, unidades y umbrales que permitan comprobar su cumplimiento.

% La tabla se presenta en una página horizontal para que sus siete columnas sean legibles.
\begin{landscape}
\begingroup
\footnotesize
\setlength{\tw}{\dimexpr\linewidth-14\tabcolsep-8\arrayrulewidth\relax}
\begin{longtable}{|L{0.05\tw}|L{0.135\tw}|L{0.195\tw}|L{0.26\tw}|L{0.13\tw}|L{0.09\tw}|L{0.14\tw}|}
\caption{Identificación y reformulación de requisitos de calidad.}\label{tab:calidad}\\
\hline
\filaencabezado \encabezado{Código} & \encabezado{Requisito original} & \encabezado{Motivo por el que es vago} & \encabezado{Requisito reformulado} & \encabezado{Métrica} & \encabezado{Unidad} & \encabezado{Umbral} \\ \hline
\endfirsthead
\multicolumn{7}{@{}l}{\footnotesize\itshape Continuación de la Tabla~\thetable} \\ \hline
\filaencabezado \encabezado{Código} & \encabezado{Requisito original} & \encabezado{Motivo por el que es vago} & \encabezado{Requisito reformulado} & \encabezado{Métrica} & \encabezado{Unidad} & \encabezado{Umbral} \\ \hline
\endhead
\hline \multicolumn{7}{r}{\footnotesize\itshape Continúa en la página siguiente} \\
\endfoot
\hline
\endlastfoot

E03 & La aplicación debe ser rápida. & La palabra «rápida» es subjetiva y no establece cuánto tiempo puede tardar una operación. Tampoco indica bajo qué condiciones debe medirse el rendimiento. Por esta razón, diferentes usuarios podrían tener opiniones distintas sobre si la aplicación responde adecuadamente. & La plataforma deberá responder a las operaciones de consulta de disponibilidad, visualización de horarios y reserva de citas en un máximo de 2 segundos para al menos el 95~\% de las solicitudes, bajo una carga de 300 usuarios concurrentes. & Tiempo de respuesta y porcentaje de solicitudes dentro del límite. & Segundos y porcentaje. & Máximo 2 segundos para al menos el 95~\% de las solicitudes. \\ \hline
E07 & Los datos de salud de los pacientes deben ser seguros. & La palabra «seguros» no especifica qué mecanismos deben utilizarse para proteger la información ni qué situaciones se deben evitar. No establece cómo se comprobará que personas no autorizadas no puedan consultar o modificar los datos médicos. & La plataforma deberá controlar el acceso a la información médica mediante autenticación y permisos según el rol del usuario. En las pruebas de seguridad, ningún usuario sin autorización deberá poder consultar o modificar información médica. & Número de accesos no autorizados exitosos. & Cantidad de accesos. & 0 accesos no autorizados exitosos. \\ \hline
E10 & La plataforma debe ser fácil de utilizar para pacientes adultos mayores. & La expresión «fácil de utilizar» depende de la experiencia y las capacidades de cada persona. El enunciado no indica qué actividad deben completar los adultos mayores ni cuánto tiempo deberían necesitar para realizarla. & Al menos el 90~\% de los pacientes adultos mayores que participen en las pruebas deberá completar una reserva de cita sin asistencia externa y en un tiempo máximo de 3 minutos. & Porcentaje de usuarios que completan la tarea sin ayuda y tiempo de realización. & Porcentaje y minutos. & Al menos 90~\% de éxito y máximo 3 minutos. \\ \hline
E14 & La plataforma debe estar siempre disponible. & La expresión «siempre disponible» no establece un porcentaje de disponibilidad ni un periodo para medirlo. Tampoco define cómo se tratarán los mantenimientos programados o las interrupciones del servicio. & La plataforma deberá mantener una disponibilidad mensual mínima del 99,5~\%, medida mediante el monitoreo del servicio. Los periodos de mantenimiento programado y comunicados previamente no se considerarán dentro del tiempo de indisponibilidad. & Porcentaje de disponibilidad mensual. & Porcentaje. & Mínimo 99,5~\% mensual. \\ \hline
E18 & La plataforma debe integrarse sin problemas con el sistema de facturación. & La expresión «sin problemas» es imprecisa porque no establece qué información debe intercambiarse, cuánto puede tardar la comunicación ni qué cantidad de errores sería aceptable. Por esta razón, no existe una condición objetiva para verificar que la integración funciona correctamente. & La plataforma deberá intercambiar la información requerida con el sistema de facturación existente mediante sus servicios web documentados. Al menos el 99~\% de las operaciones de intercambio realizadas durante las pruebas de integración deberá completarse correctamente, sin pérdida ni duplicación de información. & Porcentaje de operaciones de intercambio completadas correctamente. & Porcentaje. & Mínimo 99~\% de operaciones exitosas. \\ \hline
\end{longtable}
\endgroup
\end{landscape}

Los valores establecidos en las reformulaciones son propuestas para el proyecto del Policlínico Ribera del Quevedo. Se definieron tomando en cuenta las necesidades del servicio, la cantidad de usuarios que podrían utilizar la plataforma y las condiciones de funcionamiento descritas en el caso.

Las métricas permiten comprobar aspectos concretos, como el tiempo de respuesta, la protección de los datos, la facilidad de uso, la disponibilidad y el intercambio de información. Los umbrales propuestos sirven como criterios de aceptación para evaluar si la plataforma cumple con los niveles esperados. Estos valores no están establecidos directamente por ISO/IEC 25010:2023; la norma se utiliza como referencia para identificar las características y subcaracterísticas de calidad correspondientes~\cite{iso25010}.

\section{Comparación de resultados con una compañera}

En esta sección se registrarán únicamente los requisitos en los que existan diferencias entre la clasificación individual y la realizada por el compañero Alvarez Rai. No es necesario incluir aquellos en los que ambas clasificaciones coincidan.

\begingroup
\small
\setlength{\tw}{\dimexpr\linewidth-8\tabcolsep-5\arrayrulewidth\relax}
\renewcommand{\arraystretch}{2.2}
\begin{longtable}{|L{0.12\tw}|L{0.25\tw}|L{0.25\tw}|L{0.38\tw}|}
\caption{Requisitos con diferencias entre la clasificación propia y la de la compañera.}\label{tab:comparacion}\\
\hline
\filaencabezado \encabezado{Requisito} & \encabezado{Clasificación propia} & \encabezado{Clasificación del compañero Alvarez Rai} & \encabezado{Observación} \\ \hline
\endfirsthead
\hline
\filaencabezado \encabezado{Requisito} & \encabezado{Clasificación propia} & \encabezado{Clasificación de la compañera} & \encabezado{Observación} \\ \hline
\endhead
E10&Calidad: capacidad de interacción – inclusividad y operabilidad. &Calidad: capacidad de interacción – learnability y operabilidad. &Se coincide en la subcaracterística de operabilidad, pero existe una diferencia en la segunda subcaracterística. Se considera la inclusividad por las necesidades de los pacientes adultos mayores, mientras que la otra clasificación contempla la facilidad de aprendizaje (learnability). \\ \hline
E19&Calidad: fiabilidad – recuperabilidad. &Calidad: fiabilidad – recuperabilidad y tolerancia a fallos. &Se coincide en la recuperabilidad, pero existe una diferencia respecto a la inclusión de la tolerancia a fallos. Esta última también puede relacionarse con el enunciado, ya que contempla el fallo del servidor principal y la continuidad del servicio. \\ \hline
\end{longtable}
\endgroup

\subsection{Resumen de la comparación}

Al comparar ambas clasificaciones, se encontraron diferencias únicamente en los requisitos E10 y E19. En el E10, la diferencia corresponde a la subcaracterística de capacidad de interacción considerada; en el E19, se relaciona con la inclusión de la tolerancia a fallos junto con la recuperabilidad. En los otros 18 requisitos existe coincidencia en las clasificaciones realizadas.

\section{Evidencias del uso de herramientas de apoyo}

En este apartado se incluirán capturas de pantalla que permitan identificar las consultas realizadas a las herramientas de apoyo durante el desarrollo de la actividad.

Las evidencias se presentarán mediante figuras numeradas, cada una con un título que describa el contenido de la captura. En el texto se hará referencia a las figuras correspondientes: la Figura~\ref{fig:clasificacion} corresponde a la consulta sobre la clasificación de requisitos, la Figura~\ref{fig:reformulacion} a la consulta sobre la reformulación de requisitos de calidad y la Figura~\ref{fig:iso25010} a la consulta sobre las características de calidad de ISO/IEC 25010:2023.

\begin{figure}[H]
  \centering
  \capturaimagen{captura1}{Captura de pantalla 1\\[0.4em]Insertar aquí la captura real de la consulta relacionada con la clasificación de los 20 requisitos.\\[0.4em]Archivo esperado: imagenes/captura1.png}
  \caption{Consulta de apoyo sobre la redacción, ortografía y gramática del trabajo.}
  \label{fig:clasificacion}
\end{figure}


\Needspace{9\baselineskip}
\subsection{Registro de herramientas utilizadas}

\begingroup
\small
\setlength{\tw}{\dimexpr\linewidth-8\tabcolsep-5\arrayrulewidth\relax}
\begin{longtable}{|L{0.17\tw}|L{0.10\tw}|L{0.38\tw}|L{0.35\tw}|}
\caption{Registro de herramientas utilizadas.}\label{tab:herramientas}\\
\hline
\filaencabezado \encabezado{Herramienta} & \encabezado{Versión} & \encabezado{Propósito de uso} & \encabezado{Verificación} \\ \hline
\endfirsthead
\hline
\filaencabezado \encabezado{Herramienta} & \encabezado{Versión} & \encabezado{Propósito de uso} & \encabezado{Verificación} \\ \hline
\endhead
ChatGPT & GPT-6 & Solicitar una revisión exhaustiva de la redacción, ortografía, gramática y coherencia del trabajo académico, manteniendo su estructura y contenido original. & Se revisaron y corrigieron errores de escritura, concordancia y claridad. Se conservaron las clasificaciones, métricas, referencias y conclusiones, y se señalaron los aspectos técnicos que requieren validación. \\ \hline
\end{longtable}
\endgroup


% ---------- Referencias ----------
\printbibliography[heading=bibnumbered,title={Referencias bibliográficas}]

\end{document}
